\chapter*{Wstęp}
\addcontentsline{toc}{chapter}{Wstęp}

Dzisiejsze rozwiązania technologiczne pozwalają na łatwe zbieranie danych z czujników i przesyłanie ich do sieci w czasie rzeczywistym. Niniejszy projekt polega na stworzeniu systemu do monitorowania temperatury, który składa się z mikrokontrolera ESP32, serwera Spring Boot oraz bazy danych MySQL. System został zaprojektowany tak, aby dane \\ z urządzenia pomiarowego były zapisywane na serwerze i udostępniane użytkownikowi przez interfejs REST API.

Użytkownik ma dostęp do zebranych informacji za pomocą dwóch niezależnych interfejsów: strony internetowej oraz aplikacji mobilnej na system Android. Obie te aplikacje pobierają dane z tego samego źródła, co pozwala na bieżący podgląd temperatury oraz przeglądanie historii zapisów z datą i godziną. Projekt pokazuje, jak połączyć urządzenia IoT \\ z aplikacjami webowymi i mobilnymi w ramach jednego, spójnego systemu komunikacji.
